% Langes Solenoid im Querschnitt: B nur innen, A_phi wächst innen linear und fällt außen wie 1/r
\input{style}
\begin{document}
\begin{tikzpicture}
  \def\R{1.0}
  \fill[soft=blue] (0,0) circle (\R);
  % A_phi tangential (gegen den Uhrzeigersinn), mittig auf den Kreisen, Länge proportional zu A_phi(r)
  \foreach \r/\L in {0.55/0.3, 1.5/0.38, 2.2/0.26} {
    \draw[aux] (0,0) circle (\r);
    \foreach \w in {0,45,...,315}
      \draw[acc, vec, line width=1pt] ($(\w:\r)+({\w-90}:\L/2)$) -- ++({\w+90}:\L);
  }
  % Flächenstrom K = nI gegen den Uhrzeigersinn
  \draw[line width=1.6pt] (0,0) circle (\R);
  \draw[line width=1.6pt, -{Stealth[length=6pt,width=5pt,bend]}] (100:\R) arc[start angle=100, end angle=125, radius=\R];
  \draw[blue, line width=0.7pt] (0,0) circle (0.09); \fill[blue] (0,0) circle (0.03);
  \node[blue, font=\footnotesize, anchor=west] at (0.1,0.0) {$\vv B$};
  \node[anchor=north west, font=\footnotesize, lbl] at (-60:\R+0.05) {$K=nI$};
  \node[acc, font=\footnotesize, anchor=south] at (0,2.45) {$\vv A=A_\varphi\,\hat{\boldsymbol\varphi}$};

  % Verlauf A_phi(r)
  \begin{scope}[xshift=3.7cm, yshift=-1.0cm]
    \draw[axis] (0,0) -- (3.6,0) node[anchor=north, text=fg] {$r$};
    \draw[axis] (0,0) -- (0,2.2) node[anchor=east, text=fg] {$A_\varphi$};
    \draw[acc, line width=1.4pt] (0,0) -- (1,1.4);
    \draw[acc, line width=1.4pt, domain=1:3.4, samples=50] plot (\x,{1.4/\x});
    \draw[aux] (1,0) -- (1,1.4);
    \node[anchor=north, font=\footnotesize] at (1,0) {$R$};
    \node[acc, font=\footnotesize, anchor=east] at (-0.05,0.9) {$\tfrac12\mu_0nI\,r$};
    \node[acc, font=\footnotesize, anchor=south west] at (1.7,0.85) {$\tfrac12\mu_0nI\,\dfrac{R^2}{r}$};
  \end{scope}
\end{tikzpicture}
\end{document}
